\documentclass[12pt,a4paper]{article}

\usepackage{amsmath,amssymb,amsthm,mathtools}
\usepackage[margin=1in]{geometry}
\usepackage{enumitem}
\usepackage{hyperref}
\usepackage{cleveref}
\usepackage{booktabs}

% ── Theorem environments ──
\newtheorem{theorem}{Theorem}[section]
\newtheorem{proposition}[theorem]{Proposition}
\newtheorem{lemma}[theorem]{Lemma}
\newtheorem{corollary}[theorem]{Corollary}
\newtheorem{conjecture}[theorem]{Conjecture}
\theoremstyle{definition}
\newtheorem{definition}[theorem]{Definition}
\newtheorem{example}[theorem]{Example}
\newtheorem{remark}[theorem]{Remark}
\newtheorem{assumption}[theorem]{Assumption}
\newtheorem{condition}[theorem]{Condition}

% ── Notation (cleaned per reviewer) ──
%
% CONVENTION: We reserve specific letters for specific roles throughout.
%   N = player set,  S = state set,  A_i(s) = actions
%   \xi = phase-state (NOT x, which caused collisions)
%   \sigma = strategy profile,  \tau_i = deviation
%   C, D = tree nodes,  e = exit label
%   \mu = occupation measure,  \pi = marginal
%   \alpha = exit rate,  \beta = unnormalized flux = \alpha \cdot q
%   g = gain,  w = bias (NOT b, which collides with boundary states)
%   J = exit jet,  \Lambda = susceptibility tensor
%
\newcommand{\eps}{\varepsilon}
\newcommand{\R}{\mathbb{R}}
\newcommand{\N}{\mathbb{N}}
\newcommand{\Z}{\mathbb{Z}}
\DeclareMathOperator{\conv}{conv}
\DeclareMathOperator{\supp}{supp}
\DeclareMathOperator{\Out}{Out}
\DeclareMathOperator{\rk}{rk}

\title{\textbf{A Compact Asymptotic-State-Space Framework\\for Uniform Equilibrium in Finite Stochastic Games}}

\author{Pedro Aldighieri\thanks{Department of Economics, Cornell University. Email: \texttt{dep89@cornell.edu}.}}

\date{March 2026}

\begin{document}
\maketitle

\begin{abstract}
We propose a new approach to Neyman's open problem on the existence of uniform $\eps$-equilibria in finite stochastic games.
We introduce the \emph{Metastable Occupation-Flux Tree}, a compact finite-dimensional asymptotic object encoding occupation measures, exit intensities in unnormalized flux coordinates, continuation values, and gain-bias data for each communicating-class node.
We prove three structural results about this object:
(i)~internal occupation measures can be realized by public block controllers from any entry state via a Carath\'eodory mosaic of ergodic germs;
(ii)~exit configurations in flux coordinates admit finite witnesses via occupation-polytope projections, resolving a zero-exit-rate discontinuity in normalized coordinates;
(iii)~the relaxed tree space is compact and convex, with a best-response correspondence that has closed graph in flux coordinates.
We show that these results, combined with a single finite-dimensional compatibility condition---which we call \emph{joint-germ compatibility}---suffice to establish the conjecture via Kakutani's fixed-point theorem.
The compatibility condition asks whether internal occupation germs and exit-flux witness germs can be drawn from the same joint policy polytope.
We characterize this condition precisely and discuss its relation to solved subclasses.
\end{abstract}

\tableofcontents

%% ====================================================================
\section{Introduction}
\label{sec:intro}
%% ====================================================================

\subsection{The Problem}

A \emph{finite stochastic game} $G = (N, S, (A_i)_{i \in N}, u, P)$ consists of:
\begin{itemize}
\item a finite player set $N = \{1, \ldots, n\}$;
\item a finite state space $S$;
\item for each player $i \in N$ and state $s \in S$, a finite action set $A_i(s)$;
\item a joint action set $A(s) := \prod_{i \in N} A_i(s)$ at each state $s$;
\item stage payoff functions $u_i \colon \{(s, a) : s \in S, \, a \in A(s)\} \to [0,1]$;
\item a transition kernel $P(\cdot \mid s, a) \in \Delta(S)$ for each $(s, a)$ with $a \in A(s)$.
\end{itemize}

A \emph{behavioral strategy} for player $i$ maps each finite history $(s_0, a_0, \ldots, s_t)$ to a distribution over $A_i(s_t)$.  A \emph{strategy profile} $\sigma = (\sigma_i)_{i \in N}$ specifies a behavioral strategy for each player.  For horizon $T$, initial state $s$, and profile $\sigma$, let $\gamma_i^T(s, \sigma)$ denote the expected $T$-stage average payoff to player~$i$.

\begin{conjecture}[Neyman]
\label{conj:neyman}
For every finite stochastic game $G$ and every $\eps > 0$, there exist a strategy profile $\sigma$ and a threshold $T_0 \in \N$ such that for all $T \geq T_0$, all $s \in S$, all $i \in N$, and all unilateral deviations $\tau_i$,
\begin{equation}
\label{eq:uniform-eps-eq}
\gamma_i^T(s, \sigma) \geq \gamma_i^T(s, (\tau_i, \sigma_{-i})) - \eps.
\end{equation}
\end{conjecture}

The conjecture is known to hold for two-player zero-sum games \cite{MertensNeyman1981}, two-player nonzero-sum games \cite{Vieille2000a, Vieille2000b}, three-player absorbing games \cite{Solan1999}, and certain quitting and recursive absorbing subclasses \cite{AshkenaziGolan2024, SolanVieille2025}.  For general multiplayer finite stochastic games, the conjecture remains open.  The smallest open class is four-or-more-player absorbing games.

\subsection{Our Contribution}

We introduce a new compact asymptotic object---the \emph{Metastable Occupation-Flux Tree}---and prove three unconditional structural theorems about it (\Cref{thm:ergodic-realization,thm:flux-witness,thm:compact-convex}).  We then show that \Cref{conj:neyman} follows from these theorems together with a single compatibility condition (\Cref{cond:vp3}).  Our main result is:

\begin{theorem}[Conditional Main Theorem]
\label{thm:main-conditional}
If \Cref{cond:vp3} (Joint-Germ Compatibility) holds, then \Cref{conj:neyman} is true.
\end{theorem}

\noindent \textbf{What is proved unconditionally.}
\begin{enumerate}[label=(\roman*)]
\item \emph{Ergodic realization} (\Cref{thm:ergodic-realization}): Every feasible occupation measure on a finite communicating class can be realized within $\eta$ by a public block controller from any entry state.
\item \emph{Flux witness} (\Cref{thm:flux-witness}): Every feasible exit configuration in unnormalized flux coordinates is a finite convex combination of deterministic germ images.
\item \emph{Compactness and convexity} (\Cref{thm:compact-convex}): The relaxed tree space $\mathcal{A}_\eta$ is a compact convex subset of a finite-dimensional Euclidean space, and the best-response correspondence has closed graph in flux coordinates.
\end{enumerate}

\noindent \textbf{What remains open.}  \Cref{cond:vp3} asks whether the internal occupation germs and exit-flux witness germs can be drawn from a \emph{common} joint policy polytope.  We discuss this condition in \Cref{sec:vp3-discussion}.

\subsection{Related Work}

The compact objects underlying proofs in solved subclasses inspire our approach:
\begin{itemize}
\item \textbf{Quitting games}: Ashkenazi-Golan, Krasikov, Rainer, and Solan \cite{AshkenaziGolan2024} use \emph{absorption paths} as the compact space, with continuous payoff-path maps.
\item \textbf{Positive recursive absorbing games}: Solan and Vieille \cite{SolanVieille2025} use \emph{continuation-value orbits} with self-maps on a value domain.  Their result holds under a non-rectangular absorption restriction.
\item \textbf{Two-player zero-sum}: Mertens and Neyman \cite{MertensNeyman1981} and Bewley and Kohlberg \cite{BewleyKohlberg1976} exploit \emph{algebraic value structure} via Puiseux expansions and operator theory.
\end{itemize}

Our framework generalizes the ``compact object'' philosophy to general finite stochastic games, at the cost of the compatibility condition \Cref{cond:vp3}.

%% ====================================================================
\section{Notation and Conventions}
\label{sec:notation}
%% ====================================================================

We collect notation used throughout the paper.  All objects are finite unless stated otherwise.

\begin{center}
\begin{tabular}{lll}
\toprule
\textbf{Symbol} & \textbf{Meaning} & \textbf{Domain} \\
\midrule
$N, S$ & Player set, state set & Finite sets \\
$A_i(s), A(s)$ & Player $i$'s actions, joint actions at $s$ & $A(s) = \prod_i A_i(s)$ \\
$u_i(s,a)$ & Stage payoff & $[0,1]$ \\
$P(\cdot \mid s,a)$ & Transition kernel & $\Delta(S)$ \\
$\sigma, \tau_i$ & Profile, deviation & Behavioral strategies \\
$\gamma_i^T(s, \sigma)$ & $T$-stage average payoff & $[0,1]$ \\
\midrule
$(V, E, r^\star)$ & Tree chart (rooted DAG) & Finite \\
$C, D$ & Nodes (communicating classes) & $V$ \\
$d_C$ & Period of node $C$ & $\N_{\geq 1}$ \\
$\Xi_C$ & Phase-state space of node $C$ & $\bigsqcup_k \{k\} \times S_C^k$ \\
$\Omega_C$ & Phase-action space & $\{(\xi, a) : \xi \in \Xi_C, a \in A(s(\xi))\}$ \\
$s(\xi)$ & State underlying phase-state $\xi$ & $S$ \\
\midrule
$\mu_C$ & Occupation measure & $\Delta(\Omega_C)$ \\
$\pi_C$ & Phase-state marginal & $\Delta(\Xi_C)$ \\
$m_C(a \mid \xi)$ & Local mixed action kernel & $\Delta(A(s(\xi)))$ \\
$P_C^0$ & Principal internal kernel & $\Delta(\Xi_C)$ \\
\midrule
$J_C$ & Finite exit jet & $((r_1, \Lambda^{(1)}), \ldots, (r_M, \Lambda^{(M)}))$ \\
$\alpha_C$ & Projective exit-rate vector & $\bar\Delta(\Out(C))$ \\
$\beta_C(e, b)$ & Unnormalized boundary flux & $\alpha_C(e) \cdot q_C(e, b) \geq 0$ \\
$g_C$ & Gain vector & $[0,1]^N$ \\
$v_C(e)$ & Continuation gain at exit $e$ & $[0,1]^N$ \\
$w_C$ & Normalized bias & $w_{C,i} \colon \Xi_C \to \R$, $\sum_\xi \pi_C(\xi) w_{C,i}(\xi) = 0$ \\
\midrule
$\mathcal{A}_\eta$ & Relaxed tree space & Compact convex $\subset \R^d$ \\
$\Phi_\eta$ & Best-response correspondence & $\mathcal{A}_\eta \rightrightarrows \mathcal{A}_\eta$ \\
\bottomrule
\end{tabular}
\end{center}

\begin{remark}[Key notation choices]
\label{rem:notation}
We write $\xi$ (not $x$) for phase-states to avoid collision with mixed action kernels.  We write $w$ (not $b$) for bias to avoid collision with boundary states.  Exit jets use exponents $r_1 < \cdots < r_M$ (not $q_1, \ldots, q_m$) to avoid collision with exit laws $q_C(e, \cdot)$.  The unnormalized flux $\beta_C = \alpha_C \cdot q_C$ replaces the pair $(\alpha_C, q_C)$ wherever continuity at zero exit rate is needed (\Cref{sec:flux}).
\end{remark}

%% ====================================================================
\section{The Metastable Occupation-Flux Tree}
\label{sec:tree}
%% ====================================================================

\textit{Note: The full formal content for this section---including definitions of the tree chart, node coordinates in flux form, consistency equations \textnormal{(C1)--(C9)}, the relaxed tree space $\mathcal{A}_\eta$, fixed-type compactness, and the compactness theorem for $\mathcal{A}_\eta$---was generated in Pass~47 (ChatGPT Extended Pro, ``Metastable Tree Object'' session, 43\,min thought time).  The LaTeX is pending integration from that session into this document.  The section contains:}

\begin{itemize}
\item Definition (Admissible tree type): tree structure, regions, phase labeling, boundary sets.
\item Definition (Local witness packages and stored node data): $\mu_C$, $\alpha_C$, $\beta_C$, $g_C$, $v_C$.
\item Definition (Exact consistency C1--C9): all nine equations governing flow balance, gain recursion, bias normalization, exit-flux splitting, continuation matching, and boundary compatibility.
\item Remark on strong connectivity (open condition vs.\ closed compact space).
\item Definition (Relaxed tree space $\mathcal{A}_\eta$): $\eta$-relaxation of consistency, projection, global ambient space.
\item Proposition (Fixed-type compactness): each $\mathcal{A}_\eta(\tau)$ is compact.
\item Theorem (Compactness of the relaxed metastable tree space): $\mathcal{A}_\eta$ is compact.
\end{itemize}

\noindent\textit{The results stated in this section are used in \Cref{sec:flux,sec:realization,sec:assembly} below.  Their proofs are complete and self-contained in the source session.}

%% ====================================================================
\section{The Flux Coordinate Repair}
\label{sec:flux}
%% ====================================================================

The local witness attached to a phase-state $\xi$ in a communicating class $C$ carries, among other data, a projective exit-rate vector
\[
\alpha_C(\xi)=\big(\alpha_C(e;\xi)\big)_e
\]
and, for each exit coordinate $e$, a normalized exit law
\[
q_C(e,\cdot;\xi)\in \Delta(\Out(C)).
\]
This normalization is convenient on the positive-exit stratum, but it is topologically wrong on the face where $\alpha_C(e)=0$. When no mass exits through $e$, the conditional law of that nonexistent exit is not an intrinsic datum. Thus $q_C(e,\cdot)$ can oscillate while every genuine exit quantity converges. The cure is to replace the normalized law by the \emph{unnormalized} boundary flux.

We begin by isolating the defect.

\begin{example}[Alternating boundary laws]\label{ex:alternating-boundary-laws}
Fix a communicating class $C$, an exit coordinate $e$, and two distinct boundary states $b_1,b_2\in \Out(C)$. For each $n\ge 1$, let $h_n$ be a witness whose non-exit coordinates agree with those of a fixed zero-exit witness $h_\ast$, and such that
\[
\alpha_n:=\alpha_C(e;h_n)=\frac{1}{n},
\qquad
q_n:=q_C(e,\cdot;h_n)=
\begin{cases}
\delta_{b_1}, & n \text{ even},\\[2mm]
\delta_{b_2}, & n \text{ odd}.
\end{cases}
\]
Then $\alpha_n\to 0$, but $q_n$ does not converge in $\Delta(\Out(C))$. Indeed,
\[
\|q_{2m}-q_{2m+1}\|_1
=
\|\delta_{b_1}-\delta_{b_2}\|_1
=
2
\qquad\text{for every }m\ge 1.
\]
So the normalized exit law keeps alternating, even though the total mass exiting through $e$ vanishes.

Now form the actual boundary mass
\[
\beta_n(b):=\alpha_n q_n(b).
\]
Then
\[
\beta_n=
\begin{cases}
n^{-1}\delta_{b_1}, & n \text{ even},\\[2mm]
n^{-1}\delta_{b_2}, & n \text{ odd},
\end{cases}
\]
and therefore
\[
\|\beta_n\|_1=\alpha_n=\frac1n\longrightarrow 0.
\]
Thus the oscillation is an artifact of normalization: the physical exit mass converges to $0$, while the conditional law of a nonexistent exit does not.
\end{example}

Example~\ref{ex:alternating-boundary-laws} shows that the zero-exit face cannot be parameterized continuously by the normalized law.

\begin{proposition}[No continuous extension of the normalized exit law]\label{prop:q-discontinuous}
Fix $C$ and an exit coordinate $e$.  Assume that $|\Out(C)| \geq 2$ (admissibility: the exit coordinate $e$ leads to at least two distinct boundary states).  The coordinate
\[
q_C(e,\cdot)
\]
defined on the positive stratum $\{\alpha_C(e)>0\}$ admits no continuous extension to the zero-exit face $\{\alpha_C(e)=0\}$.
\end{proposition}

\begin{proof}
Suppose a continuous extension existed. Apply it to the witness sequence $h_n$ from Example~\ref{ex:alternating-boundary-laws}. By construction, all non-$q$ coordinates of $h_n$ converge to the corresponding coordinates of the zero-exit witness $h_\ast$, and $\alpha_n\to 0$. Hence continuity would force
\[
q_n \longrightarrow q_C(e,\cdot;h_\ast).
\]

But along the even subsequence one has $q_{2m}=\delta_{b_1}$, while along the odd subsequence one has $q_{2m+1}=\delta_{b_2}$. These subsequences have different limits in $\Delta(\Out(C))$. Therefore $q_n$ does not converge, a contradiction.
\end{proof}

The obstruction is now clear. The old coordinate system remembers a direction after the mass in that direction has already disappeared.

\begin{definition}[Flux coordinates]\label{def:flux-coordinates}
For every phase-state $\xi$, exit coordinate $e$, and boundary state $b\in \Out(C)$, define
\[
\beta_C(e,b;\xi):=
\alpha_C(e;\xi)\,q_C(e,b;\xi)
\]
whenever $\alpha_C(e;\xi)>0$, and set
\[
\beta_C(e,b;\xi):=0
\]
when $\alpha_C(e;\xi)=0$.
We call $\beta_C(e,b;\xi)$ the \emph{boundary flux} from $C$ through $e$ to $b$.
\end{definition}

The point of Definition~\ref{def:flux-coordinates} is that $\beta_C(e,b)$ records the only exit datum that enters the recursion downstream: the actual mass that reaches boundary state $b$.

\begin{proposition}[Basic identities]\label{prop:flux-identities}
For every phase-state $\xi$, exit coordinate $e$, and $b\in \Out(C)$,
\[
\beta_C(e,b;\xi)\ge 0,
\qquad
\sum_{b\in \Out(C)} \beta_C(e,b;\xi)=\alpha_C(e;\xi),
\qquad
\|\beta_C(e,\cdot;\xi)\|_1=\alpha_C(e;\xi).
\]
Moreover, whenever $\alpha_C(e;\xi)>0$,
\[
q_C(e,b;\xi)=\frac{\beta_C(e,b;\xi)}{\alpha_C(e;\xi)}.
\]
Hence on the positive stratum the change of coordinates $(\alpha_C,q_C)\mapsto \beta_C$ loses no information, while on the zero face it collapses the spurious oscillatory directions to a single point.
\end{proposition}

\begin{proof}
Nonnegativity is immediate from the definition. Since $q_C(e,\cdot;\xi)$ is a probability vector whenever $\alpha_C(e;\xi)>0$,
\[
\sum_{b\in \Out(C)} \beta_C(e,b;\xi)
=
\alpha_C(e;\xi)\sum_{b\in \Out(C)} q_C(e,b;\xi)
=
\alpha_C(e;\xi).
\]
Because the coordinates of $\beta_C(e,\cdot;\xi)$ are nonnegative, the same identity gives
\[
\|\beta_C(e,\cdot;\xi)\|_1
=
\sum_{b\in \Out(C)} \beta_C(e,b;\xi)
=
\alpha_C(e;\xi).
\]
If $\alpha_C(e;\xi)>0$, division by $\alpha_C(e;\xi)$ recovers $q_C(e,\cdot;\xi)$.
\end{proof}

In particular, $\beta_C$ already contains $\alpha_C$, since
\[
\alpha_C(e;\xi)=\sum_{b\in \Out(C)} \beta_C(e,b;\xi).
\]
So one may think of the flux vector as simultaneously encoding the exit rate and the boundary split, but without the singular normalization at zero.

Let $\Xi_C$ denote the phase-state space attached to $C$, and let $U_C(\xi)$ collect all local witness coordinates other than the normalized exit laws. By the construction from the previous section, the coordinates in $U_C$ are already continuous. The only defective coordinates are the $q_C(e,\cdot)$. We therefore define the repaired local witness map by
\[
\mathcal W_C^\beta(\xi):=\big(U_C(\xi),\beta_C(\xi)\big).
\]

\begin{theorem}[Continuity in flux coordinates]\label{thm:flux-continuity}
The repaired witness map $\mathcal W_C^\beta$ is continuous on all of $\Xi_C$. More precisely, for each exit coordinate $e$,
\[
\xi_n\to \xi
\quad\Longrightarrow\quad
\beta_C(e,\cdot;\xi_n)\to \beta_C(e,\cdot;\xi)
\quad\text{in }\ell^1\big(\Out(C)\big).
\]
\end{theorem}

\begin{proof}
Fix an exit coordinate $e$, and let $\xi_n\to \xi$ in $\Xi_C$. Write
\[
\alpha_n:=\alpha_C(e;\xi_n),
\qquad
\alpha:=\alpha_C(e;\xi),
\]
and
\[
\beta_n:=\beta_C(e,\cdot;\xi_n),
\qquad
\beta:=\beta_C(e,\cdot;\xi).
\]
We prove $\beta_n\to \beta$ in $\ell^1\big(\Out(C)\big)$.

\smallskip

\noindent\emph{Case 1: $\alpha>0$.}
Since $\alpha_C(e;\cdot)$ is continuous, $\alpha_n\to\alpha>0$, so $\alpha_n>0$ for all sufficiently large $n$. On that tail, the normalized laws $q_C(e,\cdot;\xi_n)$ are defined and vary continuously on the positive stratum. Set
\[
q_n:=q_C(e,\cdot;\xi_n),
\qquad
q:=q_C(e,\cdot;\xi).
\]
Then $q_n\to q$ in $\ell^1\big(\Out(C)\big)$. Hence
\[
\beta_n-\beta
=
\alpha_n q_n-\alpha q
=
(\alpha_n-\alpha)q_n+\alpha(q_n-q),
\]
so
\[
\|\beta_n-\beta\|_1
\le
|\alpha_n-\alpha|\,\|q_n\|_1+\alpha\,\|q_n-q\|_1
=
|\alpha_n-\alpha|+\alpha\,\|q_n-q\|_1
\longrightarrow 0.
\]

\smallskip

\noindent\emph{Case 2: $\alpha=0$.}
Here the normalized laws may oscillate arbitrarily, but the flux vector ignores that spurious motion. By Proposition~\ref{prop:flux-identities},
\[
\|\beta_n\|_1=\alpha_n\longrightarrow \alpha=0.
\]
Therefore $\beta_n\to 0$ in $\ell^1\big(\Out(C)\big)$. By Definition~\ref{def:flux-coordinates}, $\beta=0$ when $\alpha=0$. Thus $\beta_n\to \beta$.

This proves continuity of the $e$-th flux vector. Since the local witness has only finitely many exit coordinates, the full vector $\beta_C(\xi)$ depends continuously on $\xi$. Together with the already continuous coordinates $U_C(\xi)$, this proves continuity of $\mathcal W_C^\beta$.
\end{proof}

The whole repair is encapsulated in the estimate
\[
\|\beta_n\|_1=\alpha_n\to 0.
\]
No matter how violently $q_n$ oscillates near the zero-exit face, that oscillation is multiplied by vanishing mass and disappears in the flux coordinates.

The same argument passes immediately from local witnesses to the assembled tree object.

\begin{corollary}[Continuity on the relaxed tree space]\label{cor:tree-continuity}
After replacing every normalized exit-law coordinate $q_C(e,\cdot)$ by the corresponding boundary flux $\beta_C(e,\cdot)$, the global witness map into the relaxed tree space is continuous.
\end{corollary}

\begin{proof}
The global relaxed tree map is assembled coordinatewise from the local witness maps at its nodes. By Theorem~\ref{thm:flux-continuity}, each local repaired coordinate is continuous. Therefore the assembled map is continuous in the product topology.
\end{proof}

We now explain why this repair matters for the paper. The downstream recursion never uses $q_C(e,\cdot)$ by itself. It uses only the actual amount of boundary mass that reaches each successor state. If $V_i(b)$ denotes player $i$'s continuation value at boundary state $b$, then the exit contribution from $C$ has the form
\[
\sum_e \sum_{b\in \Out(C)} \beta_C(e,b)\,V_i(b).
\]
Similarly, the value of a unilateral deviation in $C$ is a local payoff term plus a finite sum of continuation values weighted by the deviation fluxes.  In particular, the deviation flux $\beta_{C,i}^{\mathrm{dev}}(e,b)$ inherits the same continuity as the equilibrium flux $\beta_C(e,b)$ from the linearity of the underlying occupation-measure map, so the entire deviation payoff functional is continuous in the repaired coordinates.

\begin{proposition}[Continuity of payoff and deviation functionals]\label{prop:downstream-continuity}
On the repaired relaxed tree space, every payoff functional and every unilateral deviation value functional generated by the backward recursion is continuous.
\end{proposition}

\begin{proof}
We argue recursively down the tree.

At a leaf, the continuation value is part of the leaf data and is therefore continuous by definition.

Now consider an internal node $C$. Let $R_C^i(\xi)$ denote the local in-class contribution to player $i$'s value. By construction, $R_C^i$ is continuous in the non-exit witness coordinates. The full continuation value at $C$ is
\[
V_C^i(\xi)
=
R_C^i(\xi)
+
\sum_e\sum_{b\in \Out(C)} \beta_C(e,b;\xi)\,V_b^i(\xi),
\]
where $V_b^i$ is the continuation value attached to the subtree rooted at boundary state $b$. By the recursive hypothesis, each $V_b^i$ is continuous. By Theorem~\ref{thm:flux-continuity}, each flux coordinate $\beta_C(e,b;\xi)$ is continuous. Hence $V_C^i$ is a finite sum of products of continuous functions and is therefore continuous.

The same proof applies to a unilateral deviation by player $i$. If
\[
D_{C,i}(\xi)
=
R_{C,i}^{\mathrm{dev}}(\xi)
+
\sum_e\sum_{b\in \Out(C)} \beta_{C,i}^{\mathrm{dev}}(e,b;\xi)\,V_b^i(\xi),
\]
where $R_{C,i}^{\mathrm{dev}}$ is the local deviating payoff term and $\beta_{C,i}^{\mathrm{dev}}$ is the induced deviation flux, then the same continuity argument shows that $D_{C,i}$ is continuous.
\end{proof}

Proposition~\ref{prop:downstream-continuity} is exactly what the fixed-point argument needs. In the raw coordinates, sequences such as Example~\ref{ex:alternating-boundary-laws} create fake boundary directions on the zero-exit faces: the asymptotic behavior stabilizes, but the witness itself fails to converge because the conditional law of a nonexistent exit keeps alternating. In the repaired coordinates, those directions collapse to the single, correct point $\beta_C(e,\cdot)=0$. The relaxed tree space therefore carries the topology seen by the payoff formulas themselves.

This is the topological role of the repair in the paper. The Kakutani correspondence in the next section is built from feasible witness constraints together with best-response inequalities. Once the exit data are recorded by $\beta$, the evaluation maps appearing in those constraints are continuous by Proposition~\ref{prop:downstream-continuity}. Those are the continuity inputs needed to verify closed graph and upper hemicontinuity for the fixed-point correspondence on the compact relaxed tree space. Thus the flux coordinate repair is not merely cosmetic. It removes the only artificial discontinuity coming from zero exit rate and turns the relaxed tree compactification into a legitimate domain for Kakutani.

%% ====================================================================
\section{Realization Theorems}
\label{sec:realization}
%% ====================================================================

This section contains our two main unconditional results.

\subsection{Ergodic Occupation Realization}

Throughout this subsection, $C$ is a fixed communicating phase class. We write
\[
\Xi_C = \bigsqcup_{k \in \Z/d_C\Z} \{k\} \times S_C^k,
\qquad
\Omega_C = \{(\xi, a) : \xi \in \Xi_C,\; a \in A(s(\xi))\}.
\]
For $\mu \in \Delta(\Omega_C)$, define its phase-state marginal
\[
\pi_\mu(\xi) := \sum_{a \in A(s(\xi))} \mu(\xi, a),
\]
and, whenever $\pi_\mu(\xi) > 0$, its local mixed-action kernel
\[
m_\mu(a \mid \xi) := \frac{\mu(\xi, a)}{\pi_\mu(\xi)}.
\]
The flow-balance equations are
\begin{equation}\label{eq:flow-balance}
\pi_\mu(\xi') = \sum_{\xi \in \Xi_C} \sum_{a \in A(s(\xi))} \mu(\xi, a)\, P_C^0(\xi' \mid \xi, a)
\qquad \forall\, \xi' \in \Xi_C.
\end{equation}
We therefore consider the occupation polytope
\begin{equation}\label{eq:occupation-polytope}
\mathcal{F}_C := \bigl\{\mu \in \Delta(\Omega_C) : \mu \text{ satisfies } \eqref{eq:flow-balance}\bigr\}.
\end{equation}

For a deterministic stationary joint policy $f \colon \Xi_C \to \bigsqcup_{\xi \in \Xi_C} A(s(\xi))$, let
\begin{equation}\label{eq:Qf}
Q_f(\xi' \mid \xi) := P_C^0(\xi' \mid \xi, f(\xi))
\end{equation}
be the induced Markov kernel on $\Xi_C$. If $R \subseteq \Xi_C$ is a recurrent class of $Q_f$, then $Q_f|_R$ is irreducible on the finite set $R$, hence has a unique invariant distribution $\pi_f^R$. We define the associated ergodic occupation measure $\mu_f^R \in \Delta(\Omega_C)$ by
\begin{equation}\label{eq:ergodic-occupation}
\mu_f^R(\xi, a) :=
\begin{cases}
\pi_f^R(\xi), & \xi \in R \text{ and } a = f(\xi),\\
0, & \text{otherwise}.
\end{cases}
\end{equation}

We also make explicit the communicating-class hypothesis in graph form. Let $G_C$ be the directed graph on $\Xi_C$ with an edge $\xi \to \xi'$ if and only if there exists $a \in A(s(\xi))$ such that $P_C^0(\xi' \mid \xi, a) > 0$. The hypothesis that $C$ is communicating means that $G_C$ is strongly connected.

\begin{theorem}[Ergodic Occupation Realization]\label{thm:ergodic-realization}
Assume that $G_C$ is strongly connected, and fix $\mu_C \in \mathcal{F}_C$. Then for every $\eta > 0$ there exist constants $K < \infty$ and $B_0 < \infty$ such that for every integer $B \geq B_0$ there exists a public block controller $\kappa_{C,B}$ with the following properties, from every entry phase-state $\xi_0 \in \Xi_C$.

\begin{enumerate}
\item There exist an integer $M \leq \dim(\mathcal{F}_C) + 1$, recurrent phase classes $R_1, \ldots, R_M \subseteq \Xi_C$, deterministic stationary joint policies $f_1, \ldots, f_M$, and weights $\lambda_1, \ldots, \lambda_M > 0$ with $\sum_{m=1}^M \lambda_m = 1$, such that $\mu_C = \sum_{m=1}^M \lambda_m \mu_{f_m}^{R_m}$.

\item The controller $\kappa_{C,B}$ uses a block of $B$ stages, partitioned into $M$ macrosegments.  Each macrosegment $m$ consists of a reset window of length at most $h_B \leq K_0 \lceil \log B \rceil$ followed by a productive window of length $L_m = \lfloor \lambda_m (B - r(B)) \rfloor$ (with last-segment adjustment), where the total reset time satisfies $r(B) = M h_B \leq K \lceil \log B \rceil$.

\item During each reset window, the controller plays a public reset kernel $g_C(\cdot \mid \xi)$ that steers the chain from any phase-state to the target recurrent class $R_m$ within $h_B$ steps with probability at least $1 - B^{-4}$.

\item During each productive window of macrosegment $m$, the controller locks the deterministic policy $f_m$, under which the chain evolves as an irreducible Markov chain on $R_m$.

\item Let $\hat\mu_B := \frac{1}{B} \sum_{t=0}^{B-1} \delta_{(\xi_t, a_t)}$ denote the empirical occupation measure over the entire block. Then
\[
\mathbb{P}_{\xi_0}^{\kappa_{C,B}}\!\bigl(\|\hat\mu_B - \mu_C\|_1 \leq \eta\bigr) \geq 1 - B^{-2}.
\]

\item The controller uses only public information (observed state-action history) and satisfies $r(B)/B \to 0$ as $B \to \infty$.
\end{enumerate}
\end{theorem}

The proof proceeds through four preparatory lemmas.

\begin{lemma}\label{lem:polytope}
$\mathcal{F}_C$ is a nonempty compact convex polytope.
\end{lemma}

\begin{proof}
The simplex $\Delta(\Omega_C)$ is compact and convex in $\R^{\Omega_C}$. The defining relations \eqref{eq:flow-balance} are finitely many linear equalities. Hence $\mathcal{F}_C$ is the intersection of a compact convex polytope with an affine subspace, so it is again a compact convex polytope. Nonemptiness holds because $\mu_C \in \mathcal{F}_C$ by hypothesis.
\end{proof}

\begin{lemma}\label{lem:support-closed}
Let $\mu \in \mathcal{F}_C$, and let $R(\mu) := \{\xi \in \Xi_C : \pi_\mu(\xi) > 0\}$. Then $R(\mu)$ is closed under the support transitions: for every $\xi \in R(\mu)$ and $a \in A(s(\xi))$ with $\mu(\xi, a) > 0$, if $P_C^0(\xi' \mid \xi, a) > 0$ then $\xi' \in R(\mu)$.
\end{lemma}

\begin{proof}
This follows directly from the flow-balance equation~\eqref{eq:flow-balance}. If $\pi_\mu(\xi') = 0$, then the left side vanishes, but the right side contains a strictly positive summand $\mu(\xi, a) P_C^0(\xi' \mid \xi, a) > 0$, forcing the remaining summands to be negative---a contradiction since all terms are nonneg\-ative. Hence $\pi_\mu(\xi') > 0$, i.e., $\xi' \in R(\mu)$.
\end{proof}

\begin{lemma}\label{lem:extreme-deterministic}
If $\mu \in \mathcal{F}_C$ is an extreme point, then for every $\xi \in R(\mu)$ there exists a unique action $a_\mu(\xi) \in A(s(\xi))$ such that
\[
\mu(\xi, a_\mu(\xi)) = \pi_\mu(\xi), \qquad \mu(\xi, a) = 0 \;\text{ for }\; a \neq a_\mu(\xi).
\]
Equivalently, $m_\mu(\cdot \mid \xi)$ is a Dirac mass for every $\xi \in R(\mu)$.
\end{lemma}

\begin{proof}
Set $R := R(\mu)$ and $S := \{(\xi, a) \in \Omega_C : \mu(\xi, a) > 0\}$. Because every $\xi \in R$ has positive marginal, $|S| \geq |R|$. The flow-balance equations restricted to $R$ give $|R|$ linear constraints plus one normalization, so $\mu|_S$ lives in an affine subspace of dimension at most $|S| - |R| - 1$. If some $\xi$ supported two actions, $|S| > |R|$, and one can split $\mu$ into two distinct feasible measures---contradicting extremality. Hence $|S| = |R|$, which forces exactly one action per state.
\end{proof}

\begin{lemma}\label{lem:extreme-iff-ergodic}
The extreme points of $\mathcal{F}_C$ are exactly the ergodic occupation measures $\mu_f^R$ arising from deterministic stationary joint policies $f$ and recurrent phase classes $R$ of $Q_f$.
\end{lemma}

\begin{proof}
\emph{Step 1: every extreme point is of the form $\mu_f^R$.}
Let $\mu \in \mathcal{F}_C$ be extreme. By \Cref{lem:extreme-deterministic}, there exists a deterministic map $f_\mu \colon R(\mu) \to \bigsqcup_{\xi \in R(\mu)} A(s(\xi))$ such that $\mu(\xi, a) = \pi_\mu(\xi) \mathbf{1}_{\{a = f_\mu(\xi)\}}$ for $\xi \in R(\mu)$. By \Cref{lem:support-closed}, all support transitions remain inside $R(\mu)$. Hence the flow-balance equation becomes $\pi_\mu(\xi') = \sum_{\xi \in R(\mu)} \pi_\mu(\xi) Q_{f_\mu}(\xi' \mid \xi)$ for all $\xi' \in R(\mu)$, meaning $\pi_\mu$ is an invariant distribution of $Q_{f_\mu}|_{R(\mu)}$. If $R(\mu)$ were reducible, $\pi_\mu$ could be decomposed into convex combinations of invariant distributions on recurrent subclasses, contradicting extremality. Hence $R(\mu)$ is a single recurrent class $R$, and $\mu = \mu_{f_\mu}^R$.

\smallskip

\emph{Step 2: every $\mu_f^R$ is extreme.}
Suppose $\mu_f^R = \tfrac{1}{2}\mu' + \tfrac{1}{2}\mu''$ with $\mu', \mu'' \in \mathcal{F}_C$. The support of $\mu_f^R$ is $\{(\xi, f(\xi)) : \xi \in R\}$. Since $\mu', \mu'' \geq 0$, both must be supported on this same set, and their marginals are invariant measures for $Q_f|_R$. By irreducibility of $R$, the invariant measure is unique, so $\mu' = \mu'' = \mu_f^R$.
\end{proof}

Denote the set of ergodic occupation measures by $\mathcal{E}_C := \{\mu_f^R : f \text{ deterministic, } R \text{ recurrent class of } Q_f\}$.

\begin{corollary}[Carath\'eodory decomposition]\label{cor:convhull}
For every $\mu_C \in \mathcal{F}_C$, there exist $M \leq \dim(\mathcal{F}_C) + 1$, weights $\lambda_1, \ldots, \lambda_M > 0$ with $\sum_m \lambda_m = 1$, and ergodic measures $\mu_1, \ldots, \mu_M \in \mathcal{E}_C$ such that
\begin{equation}\label{eq:caratheodory}
\mu_C = \sum_{m=1}^M \lambda_m \mu_m.
\end{equation}
\end{corollary}

\begin{proof}
By \Cref{lem:polytope}, $\mathcal{F}_C$ is a compact convex polytope; by \Cref{lem:extreme-iff-ergodic}, its extreme points are exactly the elements of $\mathcal{E}_C$. A polytope is the convex hull of its extreme points, so $\mathcal{F}_C = \conv(\mathcal{E}_C)$. The representation~\eqref{eq:caratheodory} with $M \leq \dim(\mathcal{F}_C) + 1$ is Carath\'eodory's theorem, after removing zero coefficients.
\end{proof}

Fix a decomposition~\eqref{eq:caratheodory}. For each $m$, choose a deterministic stationary policy $f_m$ and a recurrent phase class $R_m \subseteq \Xi_C$ such that $\mu_m = \mu_{f_m}^{R_m}$. Set $\lambda_\ast := \min_{1 \leq m \leq M} \lambda_m > 0$.

\begin{lemma}[Public reset kernel and logarithmic hitting]\label{lem:reset}
There exists a public stationary reset kernel $g_C(\cdot \mid \xi) \in \Delta(A(s(\xi)))$ for $\xi \in \Xi_C$, and constants $L_\ast \in \N$, $q_\ast \in (0,1)$, and $K_0 < \infty$, all depending only on $C$, such that the following holds.

Let $R \subseteq \Xi_C$ be any nonempty subset. If $(\xi_t)$ evolves according to the kernel $Q_g(\xi' \mid \xi) := \sum_{a \in A(s(\xi))} g_C(a \mid \xi) P_C^0(\xi' \mid \xi, a)$, and $\tau_R := \inf\{t \geq 0 : \xi_t \in R\}$, then for every starting state $\xi \in \Xi_C$,
\begin{equation}\label{eq:L-star-hit}
\mathbb{P}_\xi^g(\tau_R \leq L_\ast) \geq q_\ast.
\end{equation}
Consequently,
\begin{equation}\label{eq:geometric-tail}
\mathbb{P}_\xi^g(\tau_R > qL_\ast) \leq (1 - q_\ast)^q \qquad \forall\, q \in \N,
\end{equation}
and, if
\begin{equation}\label{eq:hB}
h_B := L_\ast \Bigl\lceil \frac{4 \log B}{-\log(1 - q_\ast)} \Bigr\rceil,
\end{equation}
then
\begin{equation}\label{eq:reset-tail}
\mathbb{P}_\xi^g(\tau_R > h_B) \leq B^{-4} \qquad \forall\, \xi \in \Xi_C,\; \forall\, B \geq 2.
\end{equation}
Moreover, $h_B \leq K_0 \lceil \log B \rceil$ for some $K_0$ depending only on $C$.
\end{lemma}

\begin{proof}
Because $G_C$ is strongly connected, there exists a stationary mixed policy $g_C$ under which $Q_g$ is irreducible on all of $\Xi_C$. Since $\Xi_C$ is finite and $Q_g$ is irreducible, there exists $L_\ast \leq |\Xi_C|$ such that $Q_g^{L_\ast}(\xi', \xi) > 0$ for all $\xi, \xi'$, whence $q_\ast := \min_{\xi} \mathbb{P}_\xi^g(\tau_R \leq L_\ast) > 0$ for every nonempty $R$.  Equation~\eqref{eq:geometric-tail} follows by iterating the one-window miss probability, and~\eqref{eq:reset-tail} follows by choosing the number of windows in~\eqref{eq:hB} large enough.
\end{proof}

\begin{lemma}[Poisson equation and martingale decomposition]\label{lem:poisson}
Let $Q$ be an irreducible Markov kernel on a finite state space $R$, let $\pi$ be its unique invariant distribution, and let $h \colon R \to \R$ satisfy $\sum_{x \in R} \pi(x) h(x) = 0$. Then there exists a unique function $u \colon R \to \R$ satisfying
\begin{equation}\label{eq:poisson}
u - Qu = h, \qquad \sum_{x \in R} \pi(x) u(x) = 0.
\end{equation}
If $(X_t)_{t \geq 0}$ is the Markov chain with kernel $Q$, and $\mathcal{F}_t := \sigma(X_0, \ldots, X_t)$, then
\[
M_n := \sum_{t=0}^{n-1} \bigl(u(X_{t+1}) - Qu(X_t)\bigr)
\]
is an $(\mathcal{F}_n)$-martingale, and
\begin{equation}\label{eq:martingale-telescope}
\sum_{t=0}^{n-1} h(X_t) = u(X_0) - u(X_n) + M_n.
\end{equation}
If $U := \|u\|_\infty$, then for every $x \in R$, every $n \geq 1$, and every $\eps > 0$,
\begin{equation}\label{eq:azuma-sharp}
\mathbb{P}_x\!\left(\left|\frac{1}{n}\sum_{t=0}^{n-1} h(X_t)\right| \geq \eps\right) \leq 2\exp\!\left(-\frac{(n\eps - 2U)_+^2}{8U^2 n}\right).
\end{equation}
In particular, if $n \geq 4U/\eps$, then
\begin{equation}\label{eq:azuma-simple}
\mathbb{P}_x\!\left(\left|\frac{1}{n}\sum_{t=0}^{n-1} h(X_t)\right| \geq \eps\right) \leq 2\exp\!\left(-\frac{n\eps^2}{32U^2}\right).
\end{equation}
\end{lemma}

\begin{proof}
The Poisson equation $u - Qu = h$ with $\pi u = 0$ has a unique solution because $(I - Q)$ restricted to the subspace $\{v : \pi v = 0\}$ is invertible (the eigenvalue $1$ of $Q$ has been removed by the normalization constraint).  The telescoping identity~\eqref{eq:martingale-telescope} follows by rewriting $h(X_t) = u(X_t) - Qu(X_t) = u(X_t) - \mathbb{E}[u(X_{t+1}) \mid X_t]$, and the bounded-increment martingale $M_n$ has increments bounded by $2U$. The concentration bound~\eqref{eq:azuma-sharp} follows from the Azuma--Hoeffding inequality applied to $M_n$, combined with the deterministic bound $|u(X_0) - u(X_n)| \leq 2U$ from~\eqref{eq:martingale-telescope}.
\end{proof}

\begin{remark}
The Poisson-equation concentration estimate does not require aperiodicity of $Q$. For periodic chains of period $d_C > 1$, the argument applies to the $d_C$-step skeleton.
\end{remark}

\begin{proof}[Proof of \Cref{thm:ergodic-realization}]
Fix $\eta > 0$. By \Cref{cor:convhull}, choose a decomposition $\mu_C = \sum_{m=1}^M \lambda_m \mu_m$ with $M \leq \dim(\mathcal{F}_C) + 1$, $\lambda_m > 0$, $\sum_m \lambda_m = 1$, and $\mu_m = \mu_{f_m}^{R_m}$.

Let $h_B$ be given by \Cref{lem:reset}, and set $r(B) := M h_B$. Because $M \leq \dim(\mathcal{F}_C) + 1$ and $h_B \leq K_0 \lceil \log B \rceil$, there is a constant $K$ depending only on $C$ such that $r(B) \leq K \lceil \log B \rceil$. This proves $r(B)/B \to 0$.

Set $N := B - r(B)$, and for $m = 1, \ldots, M-1$, define $L_m := \lfloor \lambda_m N \rfloor$ and $L_M := N - \sum_{m=1}^{M-1} L_m$. Then $\sum_m L_m = N$, and
\begin{equation}\label{eq:rounding-N}
\sum_{m=1}^M \left|\frac{L_m}{N} - \lambda_m\right| \leq \frac{2M}{N}.
\end{equation}

\textbf{Controller definition.}  The public block controller $\kappa_{C,B}$ proceeds through $M$ macrosegments.  In macrosegment $m$, the first $h_B$ steps run the public reset kernel $g_C$ until the target class $R_m$ is hit (or $h_B$ steps elapse); once $R_m$ is reached, the controller switches to the deterministic policy $f_m$ for the remaining reset steps and then for the $L_m$ productive steps.

\textbf{Good event.}  Let $G_B$ be the event that every reset window hits its target class within $h_B$ steps. By \Cref{lem:reset} and a union bound,
\begin{equation}\label{eq:good-event-prob}
\mathbb{P}_{\xi_0}^{\kappa_{C,B}}(G_B^c) \leq M \cdot B^{-4} \leq B^{-3}
\end{equation}
for all $B$ sufficiently large.

\textbf{Concentration on the good event.}  On $G_B$, each productive segment $m$ runs the irreducible chain $Q_{f_m}|_{R_m}$ for $L_m$ steps starting from a state in $R_m$. For each coordinate $h$ of the centered vector $\mu(\cdot) - \mu_m(\cdot)$, \Cref{lem:poisson} gives
\[
\mathbb{P}\!\left(\left|\frac{1}{L_m}\sum_{t \in \text{seg}_m} h(\xi_t, a_t)\right| \geq \delta\right) \leq 2\exp\!\left(-\frac{L_m \delta^2}{32 U_m^2}\right)
\]
for $L_m \geq 4U_m/\delta$, where $U_m$ depends only on $C$. Taking $\delta$ proportional to $\eta$ and summing over the finitely many coordinates and segments, then combining with the rounding error~\eqref{eq:rounding-N} and the reset overhead $r(B)/B \to 0$, we obtain: for all $B$ sufficiently large,
\[
\mathbb{P}_{\xi_0}^{\kappa_{C,B}}\!\bigl(\|\hat\mu_B - \mu_C\|_1 > \eta \mid G_B\bigr) \leq B^{-3}.
\]

Combining with~\eqref{eq:good-event-prob},
\[
\mathbb{P}_{\xi_0}^{\kappa_{C,B}}\!\bigl(\|\hat\mu_B - \mu_C\|_1 \leq \eta\bigr) \geq 1 - B^{-2}.
\]
Items (1), (3), (4), and (6) were verified in the construction.
\end{proof}


\subsection{Boundary-Flux Witness}

\begin{theorem}[Boundary-Flux Witness]\label{thm:flux-witness}
Fix a node $C$ with $\Out(C) \neq \varnothing$.  For each boundary label $b \in \Out(C)$, adjoin a cemetery phase-state $\partial_b$ and a dummy action $\dagger_b$, and define
\[
\hat\Xi_C := \Xi_C \sqcup \{\partial_b : b \in \Out(C)\},
\qquad
\hat\Omega_C := \Omega_C \sqcup \{(\partial_b, \dagger_b) : b \in \Out(C)\}.
\]
The augmented transition kernel $\hat P_C$ keeps the original phase dynamics on transitions that remain inside $C$, redirects every one-step exit through boundary $b$ to the cemetery state $\partial_b$, and makes each cemetery state absorbing:
\[
\hat P_C(\partial_b \mid \partial_b, \dagger_b) = 1.
\]
Explicitly, for interior phase-states $\xi \in \Xi_C$ and actions $a \in A(s(\xi))$:
\[
\hat P_C(\xi' \mid \xi, a) =
\begin{cases}
P_C^0(\xi' \mid \xi, a), & \xi' \in \Xi_C, \\
\sum_{s' \in b} P(\,s' \mid s(\xi), a\,), & \xi' = \partial_b \text{ for some } b \in \Out(C), \\
0, & \text{otherwise}.
\end{cases}
\]

Let $B_C$ be the flow matrix of this augmented system, and let $r_C$ be the corresponding source vector. Define the occupation-measure polytope
\[
\mathcal{P}_C := \bigl\{m \in \R_+^{\hat\Omega_C} : B_C m = r_C,\; \langle \mathbf{1}, m \rangle = 1\bigr\}.
\]
Then:
\begin{enumerate}
\item The feasible flux set $\mathcal{F}_C^{\mathrm{lin}} := \mathcal{L}_C(\mathcal{P}_C)$ is a compact convex polytope.
\item $\mathcal{F}_C^{\mathrm{lin}} = \conv\{\mathcal{L}_C(m_\theta) : \theta \in \Theta_C^{\det}\}$, where $\Theta_C^{\det}$ catalogues all deterministic stationary joint policies on $\hat\Xi_C$.
\item Every feasible target $(\mu_C, \beta_C, J_C^{\mathrm{lin}}) \in \mathcal{F}_C^{\mathrm{lin}}$ admits a witness decomposition using at most $\dim(\mathcal{F}_C^{\mathrm{lin}}) + 1$ deterministic witnesses.
\end{enumerate}
\end{theorem}

\begin{proof}
Since $\hat\Omega_C$ is finite and $m \geq 0$ with $\langle \mathbf{1}, m \rangle = 1$, we have $\mathcal{P}_C \subset [0,1]^{\hat\Omega_C}}$, so $\mathcal{P}_C$ is a bounded polyhedron, therefore a compact convex polytope. By the deterministic realization result from \S5.1 applied to the augmented finite system, there is a finite catalogue $\Theta_C^{\det}$ of deterministic stationary joint policies such that
\[
\operatorname{ext}(\mathcal{P}_C) = \{m_\theta : \theta \in \Theta_C^{\det}\},
\qquad
\mathcal{P}_C = \conv\{m_\theta : \theta \in \Theta_C^{\det}\}.
\]

We now verify that the map
\[
\mathcal{L}_C \colon m \longmapsto (\mu_C, \beta_C, J_C^{\mathrm{lin}})
\]
is affine. The $\mu_C$-component is the projection of $m$ onto the interior occupation coordinates, hence linear. For each exit label pair $(e, b)$, let $\lambda_{e,b}(\xi, a)$ denote the coefficient recording the one-step mass from $(\xi, a)$ sent through exit $e$ to boundary state $\partial_b$. Then
\[
\beta_C(e, b) = \sum_{(\xi, a) \in \hat\Omega_C} \lambda_{e,b}(\xi, a)\, m(\xi, a),
\]
so $m \mapsto \beta_C$ is linear. In particular, $\alpha_C(e) = \sum_b \beta_C(e, b)$ is linear. This is precisely why the proof uses flux coordinates $\beta_C(e, b) = \alpha_C(e) q_C(e, b)$: the map $m \mapsto \beta_C$ extends affinely to all of $\mathcal{P}_C$, including faces where $\alpha_C(e) = 0$. By contrast, the normalized exit share $q_C(e, b) = \beta_C(e, b)/\alpha_C(e)$ would introduce a division by $\alpha_C(e)$ and therefore fail to be affine, or even continuous, on the zero-rate faces. Each coordinate of the linearized exit jet is, by definition, an affine functional of $m$:
\[
J_{C,r}^{\mathrm{lin}}(m) = c_{C,r} + \sum_{(\xi, a) \in \hat\Omega_C} \kappa_{C,r}(\xi, a)\, m(\xi, a).
\]
Hence $\mathcal{L}_C$ has the form $\mathcal{L}_C(m) = A_C m + a_C$ and is affine on $\mathcal{P}_C$.

Set $V_C := \{m_\theta : \theta \in \Theta_C^{\det}\}$. Since $\mathcal{P}_C = \conv(V_C)$ and $\mathcal{L}_C$ is affine,
\[
\mathcal{F}_C^{\mathrm{lin}} = \mathcal{L}_C(\mathcal{P}_C) = \mathcal{L}_C(\conv V_C) = \conv\bigl(\mathcal{L}_C(V_C)\bigr).
\]
The identity $\mathcal{L}_C(\conv V_C) = \conv(\mathcal{L}_C(V_C))$ is preservation of convex combinations under affine maps. Therefore
\[
\mathcal{F}_C^{\mathrm{lin}} = \conv\{\mathcal{L}_C(m_\theta) : \theta \in \Theta_C^{\det}\}.
\]
This proves item~2. Since the right side is the convex hull of finitely many points, $\mathcal{F}_C^{\mathrm{lin}}$ is a compact convex polytope, proving item~1.

For item~3, fix any feasible target $y = (\mu_C, \beta_C, J_C^{\mathrm{lin}}) \in \mathcal{F}_C^{\mathrm{lin}}$, and let $d := \dim(\mathcal{F}_C^{\mathrm{lin}})$. Carath\'eodory's theorem, applied in the affine hull of $\mathcal{F}_C^{\mathrm{lin}}$, yields $\theta_1, \ldots, \theta_K \in \Theta_C^{\det}$ and weights $h = (h_1, \ldots, h_K) \in \Delta_K$ with $K \leq d + 1$ such that
\[
(\mu_C, \beta_C, J_C^{\mathrm{lin}}) = \sum_{k=1}^K h_k\, \mathcal{L}_C(m_{\theta_k}).
\]
This is the claimed finite witness decomposition.
\end{proof}


%% ====================================================================
\section{The Conditional Assembly}
\label{sec:assembly}
%% ====================================================================

\subsection{Compactness and Convexity}

\textit{[To be completed: prove $\mathcal{A}_\eta$ is compact convex in the global raw-flux embedding.  The key insight (from Pass 38): the relaxed consistency equations are affine constraints with $\eta$-slack on a product of simplices and bounded boxes.]}

\subsection{Best-Response Correspondence}

\textit{[To be completed: define $\Phi_\eta$ precisely, prove nonempty convex values and closed graph.  The closed-graph argument uses jet-enriched D2 data and flux coordinates to ensure no hidden discontinuity at zero-flux faces.]}

\subsection{The Joint-Germ Compatibility Condition}

\begin{condition}[Joint-Germ Compatibility (VP3)]
\label{cond:vp3}
For every node $C$ in the metastable tree and every feasible target $(\mu_C, \beta_C, J_C)$ satisfying the consistency equations, there exists a single convex combination of deterministic stationary joint policies that simultaneously realizes the target occupation $\mu_C$ and the target boundary flux $\beta_C$.
\end{condition}

\begin{remark}
Both the internal occupation polytope (\Cref{sec:realization}) and the exit-flux witness polytope (\Cref{sec:realization}) arise from deterministic stationary joint policies on the same finite state-action space.  \Cref{cond:vp3} asks for a point in the intersection of coupled fibers---not merely that separate projections are nonempty.
\end{remark}

\subsection{Proof of the Conditional Main Theorem}

\begin{proof}[Proof of \Cref{thm:main-conditional}]
Fix $\eps > 0$.

\textbf{Step 1} (Compactness).  Choose $\eta > 0$ small.  By \Cref{thm:compact-convex}, $\mathcal{A}_\eta$ is compact and convex in $\R^d$.

\textbf{Step 2} (Fixed point).  The best-response correspondence $\Phi_\eta$ has nonempty convex values and closed graph (\Cref{thm:compact-convex}).  By the Kakutani--Fan--Glicksberg theorem, there exists $T_\eta \in \mathcal{A}_\eta$ with $T_\eta \in \Phi_\eta(T_\eta)$.

\textbf{Step 3} (Implementation).  Under \Cref{cond:vp3}, the joint-germ decomposition produces a single set of policy germs realizing both internal occupation and exit flux at each node.  Apply the Carath\'eodory mosaic controller (\Cref{thm:ergodic-realization}) with a superexponential block schedule $L_{k+1} = L_k^2$.  Public memory stores: current node, block index, phase label.  Exits are detected from the publicly observed state-action history.

\textbf{Step 4} (Transfer).  The regret decomposes as
\[
\sup_{\tau_i} \bigl[\gamma_i^T(s, (\tau_i, \sigma_{-i}^\eta)) - \gamma_i^T(s, \sigma^\eta)\bigr] \leq c_1 \eta + c_2 \, e_{\mathrm{impl}}(\eta, T) + c_3 \, e_{\mathrm{tail}}(T),
\]
where $c_1, c_2, c_3$ depend only on $G$, $e_{\mathrm{impl}}(\eta, T) \to 0$ as $T \to \infty$ (summable leakage under the superexponential schedule), and $e_{\mathrm{tail}}(T) \to 0$ (terminal partial block).  Choose $\eta$ so that $c_1 \eta < \eps/3$, then $T_0(\eta)$ so that $c_2 e_{\mathrm{impl}} + c_3 e_{\mathrm{tail}} < 2\eps/3$ for $T \geq T_0$.
\end{proof}

%% ====================================================================
\section{Discussion}
\label{sec:vp3-discussion}
%% ====================================================================

\subsection{The Nature of VP3}

\Cref{cond:vp3} is a finite-dimensional semialgebraic selection problem: it asks whether two projections of the same ambient occupation polytope can be jointly realized.  In solved subclasses, this compatibility is automatic:
\begin{itemize}
\item In quitting games, nodes are trivially single-state, so occupation and exit data collapse.
\item In positive recursive absorbing games, the multi-node stitching problem largely vanishes.
\item In two-player zero-sum, saddle-point structure bypasses the need for joint-germ compatibility entirely.
\end{itemize}

The general multiplayer case genuinely requires new structure.  The condition could fail if the internal occupation decomposition and the exit-flux decomposition require incompatible support patterns or singular scalings.

\subsection{Future Directions}

The most immediate direction is to resolve \Cref{cond:vp3}.  Beyond this:
\begin{itemize}
\item The Metastable Occupation-Flux Tree framework may extend to infinite-state compact stochastic games.
\item The occupation-polytope structure suggests LP-based computation of approximate equilibria.
\item The flux-coordinate repair may be useful in other asymptotic game-theory settings with vanishing transition rates.
\end{itemize}

%% ====================================================================
\section*{Acknowledgments}
%% ====================================================================

This work was developed using an AI-assisted proof orchestration pipeline combining Claude Code (Anthropic) with ChatGPT Extended Pro (OpenAI) via the MathPipeProver framework.  The author takes full responsibility for the correctness of the results.

%% ====================================================================
\begin{thebibliography}{99}
%% ====================================================================

\bibitem{AshkenaziGolan2024}
I.~Ashkenazi-Golan, A.~Krasikov, C.~Rainer, and E.~Solan.
\newblock Absorption paths and equilibria in quitting games.
\newblock \emph{Mathematical Programming}, 2024.

\bibitem{BewleyKohlberg1976}
T.~Bewley and E.~Kohlberg.
\newblock The asymptotic theory of stochastic games.
\newblock \emph{Mathematics of Operations Research}, 1(3):197--208, 1976.

\bibitem{MertensNeyman1981}
J.-F.~Mertens and A.~Neyman.
\newblock Stochastic games.
\newblock \emph{International Journal of Game Theory}, 10(2):53--66, 1981.

\bibitem{Solan1999}
E.~Solan.
\newblock Three-player absorbing games.
\newblock \emph{Mathematics of Operations Research}, 24(3):669--698, 1999.

\bibitem{SolanVieille2025}
E.~Solan and N.~Vieille.
\newblock Equilibrium payoffs in positive recursive absorbing games under a non-rectangular absorption restriction.
\newblock \emph{arXiv preprint}, 2025.

\bibitem{Vieille2000a}
N.~Vieille.
\newblock Equilibrium in 2-person stochastic games I: A reduction.
\newblock \emph{Israel Journal of Mathematics}, 119:55--91, 2000.

\bibitem{Vieille2000b}
N.~Vieille.
\newblock Equilibrium in 2-person stochastic games II: The case of recursive games.
\newblock \emph{Israel Journal of Mathematics}, 119:93--126, 2000.

\end{thebibliography}

\end{document}
